% State-specific — Personal Financial Literacy
\section{Personal Financial Literacy}

\begin{learningGoals}
\begin{itemize}
    \item Compare financial institutions and payment methods
    \item Understand credit, debit, and how interest affects decisions
\end{itemize}
\end{learningGoals}

\begin{conceptBox}{Money Management Basics}
\begin{itemize}[leftmargin=*]
    \item \textbf{Bank} --- holds your money, offers savings and checking accounts
    \item \textbf{Credit union} --- similar to a bank, owned by its members (often lower fees)
    \item \textbf{Debit card} --- spends money you already have in your account
    \item \textbf{Credit card} --- borrows money you must pay back, often with \textbf{interest}
    \item \textbf{Interest rate} --- the percent a lender charges (or a bank pays you) on money
\end{itemize}
\end{conceptBox}

\begin{workedExample}{Comparing Payment Methods}
You buy a $\$200$ game console.

\textbf{Debit card:} You pay $\$200$ from your account. Total cost: $\$200$.

\textbf{Credit card at $18\%$ annual rate, paid after $1$ year:}

Interest $= 200 \times 0.18 \times 1 = \$36$. Total cost: $\$236$.
\answerHighlight{Credit costs $\$36$ more if you wait a year to pay.}
\end{workedExample}

\tipBox{Pay credit card bills on time every month to avoid extra interest charges!}

\begin{practiceBox}{Financial Literacy}
\resetProblems

\prob You charge $\$150$ on a credit card at $15\%$ annual interest. What is the interest after $1$ year?
\answerExplain{$\$22.50$}{$I = 150 \times 0.15 \times 1 = \$22.50$.}

\prob Bank A pays $2\%$ interest on savings. Bank B pays $3.5\%$. You deposit $\$1{,}000$ for $1$ year. How much more does Bank B earn?
\answerExplain{$\$15$ more}{Bank A: $1{,}000 \times 0.02 = \$20$. Bank B: $1{,}000 \times 0.035 = \$35$. Difference: $35 - 20 = \$15$.}

\trueOrFalse{A debit card lets you spend more than you have in your account.}
\answerTF{False}

\prob Name one advantage and one disadvantage of using a credit card.
\answerExplain{Advantage: buy now; Disadvantage: interest charges}{Credit lets you purchase before you have the cash, but you pay interest if you don't pay the full balance on time.}

\prob A checking account charges a $\$5$ monthly fee. A credit union charges $\$2$/month. How much do you save per year at the credit union?
\answerExplain{$\$36$}{Bank fee per year: $5 \times 12 = \$60$. Credit union: $2 \times 12 = \$24$. Savings: $60 - 24 = \$36$.}

\end{practiceBox}
